%! Solving Exponential Equations with a Common Base — Unit 6, Lesson 6.3 — Guided Notes
\documentclass[10pt]{article}
\usepackage{algebra2-article}
\usepackage{algebra2-boxes}
\newcommand{\TallMath}[1]{$\displaystyle #1\rule[-1.4em]{0pt}{3.2em}$}
\newcommand{\lgrid}[4]{%
  \draw[step=1, linegray!40, very thin] (#1,#3) grid (#2,#4);
  \draw[->, semithick, charcoal] (#1,0)--(#2,0) node[below right, font=\scriptsize]{$x$};
  \draw[->, semithick, charcoal] (0,#3)--(0,#4) node[left, font=\scriptsize]{$y$};}
% y = a*b^(x-h) + k, plotted as a*exp((ln b)(x-h))+k: pgfmath's pow() is unreliable for
% negative and fractional exponents, but exp() is not.
% #1=a  #2=ln(b)  #3=h  #4=k  #5=xmin  #6=xmax
% ln2 = 0.693147   ln3 = 1.098612   ln(1/2) = -0.693147
\newcommand{\expgraph}[6]{\draw[very thick, forest, domain=#5:#6, samples=140, smooth]
  plot (\x,{(#1)*exp((#2)*((\x)-(#3)))+(#4)});}
% Horizontal asymptote (navy dashed): {xmin}{xmax}{k}
\newcommand{\hasym}[3]{\draw[navy, dashed, semithick] (#1,#3)--(#2,#3);}
% The horizontal line we are solving against (gold): {xmin}{xmax}{c}
\newcommand{\solline}[3]{\draw[goldacc, very thick] (#1,#3)--(#2,#3);}

\begin{document}

\pageheader{Unit 6, Lesson 6.3}{Guided Notes}
\namedateperiod

\vspace{0.08in}

\begin{objectivebox}
\textbf{By the end of this lesson, I will be able to\dots}
\begin{itemize}\setlength{\itemsep}{2pt}
  \item solve an exponential equation by rewriting both sides as powers of \emph{one} base and then setting the exponents equal;
  \item explain \emph{why} that move is legal --- an exponential function never produces the same output twice, so matching bases forces the exponents to match;
  \item read a solution off a graph as the point where a horizontal line crosses the curve, and decide from the \emph{range} when an exponential equation has no solution at all.
\end{itemize}
\end{objectivebox}

\vspace{0.05in}

\begin{vocabbox}
\small
Fill in each term as we build it together.
\par\vspace{2pt}   % \par is required: \termblanklong's \noindent is a no-op mid-paragraph
\termblanklong{Exponential equation}
\termblanklong{Common base}
\termblanklong{One-to-one property of exponents}
\termblanklong{Power of a power, $\left(b^{m}\right)^{n}=b^{mn}$}
\termblanklong{Solution as an intersection}
\end{vocabbox}

\begin{hookbox}
% The lead-in and the two graphs must not be separated by a page break.
\begin{minipage}{\linewidth}
Two equations. Both say ``something equals $4$.'' One has \emph{two} answers and the other has exactly
\emph{one} --- and the graphs say why before you do any algebra.
\begin{center}
\begin{tikzpicture}[scale=0.38]
  % LEFT: y = x^2 meets y = 4 twice
  \begin{scope}
    \lgrid{-3}{3}{-1}{9}
    \begin{scope}
      \clip (-3,-1) rectangle (3,9);
      \draw[very thick, forest, domain=-3:3, samples=90, smooth] plot (\x,{(\x)*(\x)});
    \end{scope}
    \solline{-3}{3}{4}
    \fill[charcoal] (-2,4) circle (3pt);
    \fill[charcoal] (2,4) circle (3pt);
    \node[charcoal, font=\small] at (0,-2.6) {$y=x^{2}$ \; and \; $y=4$};
  \end{scope}
  % RIGHT: y = 2^x meets y = 4 once
  \begin{scope}[xshift=8.5cm]
    \lgrid{-3}{3}{-1}{9}
    \begin{scope}
      \clip (-3,-1) rectangle (3,9);
      \expgraph{1}{0.693147}{0}{0}{-3}{3}
    \end{scope}
    \solline{-3}{3}{4}
    \fill[charcoal] (2,4) circle (3pt);
    \node[charcoal, font=\small] at (0,-2.6) {$y=2^{x}$ \; and \; $y=4$};
  \end{scope}
\end{tikzpicture}
\end{center}
\end{minipage}
\vspace{4pt}
\begin{itemize}\setlength{\itemsep}{1pt}
  \item Solve $x^{2}=4$: \; $x=$ \blank{2.2cm} \qquad Solve $2^{x}=4$: \; $x=$ \blank{2.2cm}
  \item The gold line is at the same height in both pictures. Why does it hit one curve twice and the other only once? \writeline
  \item Which fact about the \emph{shape} of an exponential graph (Lesson 6.0) guarantees the right-hand picture can never look like the left-hand one? \writeline
\end{itemize}
Everything today rests on that last answer: \textbf{an exponential function never uses the same output
twice, so if the bases match, the exponents have no choice but to match too.}
\end{hookbox}

\begin{notesbox}{1. Why You Are Allowed to Drop the Base}
An \textbf{exponential equation} is an equation with the variable in an \blank{2.6cm}.

\vspace{4pt}
Lesson 6.0 established that an exponential graph has \textbf{no turning points} --- it is always
increasing or always decreasing, never both. A function like that can never come back to an output it
has already used, so a horizontal line crosses it \blank{2.6cm} time(s) at most.

\vspace{4pt}
\begin{center}
\begin{tcolorbox}[colback=goldbg, colframe=goldacc, boxrule=0.8pt, arc=1.5mm, width=0.90\linewidth,
    left=3mm, right=3mm, top=2mm, bottom=2mm]
\centering
\textbf{One-to-One Property of Exponents} \\[3pt]
If \; $b^{\,m} = b^{\,n}$ \; and \; $b>0$, $b\neq1$, \qquad then \; $m=n$.
\end{tcolorbox}
\end{center}

\vspace{2pt}
\textbf{The fine print is not decoration.} Lesson 6.1 told you an exponential model needs $b>0$ and
$b\neq1$. Here is what goes wrong when $b=1$: \; $1^{5}=$ \blank{1.0cm} \; and \; $1^{9}=$
\blank{1.0cm}, \; so $1^{5}=1^{9}$ even though $5\neq9$. With $b=1$ the graph is a flat line, which
repeats \emph{every} output --- matching exponents would be nonsense. The restriction you first met as
a rule about \emph{models} turns out to be exactly the condition that makes today's \emph{solving}
legal.

\vspace{3pt}
\emph{And $b>0$?} A negative base breaks at fractional exponents: $(-4)^{1/2}=\sqrt{-4}$ is not a real
number (Unit 5).
\end{notesbox}

\begin{notesbox}{2. Getting the Bases to Match}
The property applies only when the two sides are powers of the \textbf{same base}, so most of the work
in this lesson is the rewriting you did in the Warm-Up. Fill in every exponent.

\vspace{3pt}
\begin{center}
\small
\renewcommand{\arraystretch}{1.4}
\begin{tabular}{|l|c|c|c|c|c|c|}
\hline
\textbf{Powers of $2$} & $4$ & $8$ & $16$ & $32$ & $64$ & $\sqrt{2}$ \\
\textbf{exponent on $2$} & \blank{0.9cm} & \blank{0.9cm} & \blank{0.9cm} & \blank{0.9cm} & \blank{0.9cm} & \blank{0.9cm} \\
\hline
\textbf{Powers of $3$} & $9$ & $27$ & $81$ & $\frac13$ & $\frac19$ & $\frac{1}{27}$ \\
\textbf{exponent on $3$} & \blank{0.9cm} & \blank{0.9cm} & \blank{0.9cm} & \blank{0.9cm} & \blank{0.9cm} & \blank{0.9cm} \\
\hline
\textbf{Powers of $5$} & $25$ & $125$ & $\frac15$ & $\frac{1}{25}$ & $\sqrt{5}$ & $\sqrt[3]{5}$ \\
\textbf{exponent on $5$} & \blank{0.9cm} & \blank{0.9cm} & \blank{0.9cm} & \blank{0.9cm} & \blank{0.9cm} & \blank{0.9cm} \\
\hline
\end{tabular}
\end{center}

\vspace{4pt}
\begin{tcolorbox}[colback=redbg, colframe=redacc, boxrule=0.8pt, arc=1.5mm,
    left=2.5mm, right=2.5mm, top=1.5mm, bottom=1.5mm]
\textbf{The two mistakes this lesson is made of.} Both are Unit 5 exponent rules, not new ideas.
\begin{itemize}[leftmargin=1.4em, itemsep=2pt]
  \item \textbf{Power of a power multiplies.} $\left(2^{3}\right)^{x}$ is $2$ to the power
        \blank{1.6cm} --- \emph{not} $2^{\,3+x}$. You checked this numerically in the Warm-Up:
        $\left(2^{3}\right)^{4}=4096=2^{12}$, while $2^{3}\cdot2^{4}=128=2^{7}$.
  \item \textbf{Distribute across the \emph{whole} exponent.} $\left(3^{2}\right)^{x+4}=3^{\,2(x+4)}$,
        and $2(x+4)=$ \blank{2.2cm}. Writing $3^{\,2x+4}$ multiplies the $x$ and forgets the $4$.
\end{itemize}
\end{tcolorbox}
\end{notesbox}

\begin{notesbox}{3. The Method: Three Steps}
\begin{center}
\small
\renewcommand{\arraystretch}{1.3}
\begin{tabularx}{0.94\linewidth}{@{}l X@{}}
\textbf{Step 1} & Rewrite \emph{both} sides as powers of one common base. \\
\textbf{Step 2} & The bases match, so set the \blank{2.6cm} equal to each other. \\
\textbf{Step 3} & Solve what is left --- it is usually a linear equation --- and check. \\
\end{tabularx}
\end{center}

\vspace{4pt}
\begin{center}
\small
\renewcommand{\arraystretch}{1.6}
\begin{tabularx}{\linewidth}{@{}l X X c@{}}
\hline
\textbf{Equation} & \textbf{Step 1: one common base} & \textbf{Step 2: exponents equal} & \textbf{Step 3: $x=$} \\
\hline
\textbf{A.} $2^{\,x+1}=32$ & \blank{3.4cm} & \blank{3.0cm} & \blank{1.4cm} \\
\textbf{B.} $8^{x}=4$ & \blank{3.4cm} & \blank{3.0cm} & \blank{1.4cm} \\
\textbf{C.} $3^{x}=\frac{1}{27}$ & \blank{3.4cm} & \blank{3.0cm} & \blank{1.4cm} \\
\textbf{D.} $9^{\,x-2}=27^{x}$ & \blank{3.4cm} & \blank{3.0cm} & \blank{1.4cm} \\
\hline
\end{tabularx}
\end{center}

\vspace{3pt}
\textbf{Checking B is worth doing out loud.} Its answer is a fraction, which always feels wrong the
first time. Substitute it back with Lesson 5.1's rational exponents:
$8^{\,2/3}=\left(\sqrt[3]{8}\right)^{2}=$ \blank{1.6cm}, and the right side of B was
\blank{1.4cm}. \checkmark

\vspace{2pt}
A non-integer answer is not a red flag. The exponent is a \emph{location on the $x$-axis}, and nothing
requires it to land on a whole number.
\end{notesbox}

\begin{notesbox}{4. The Same Answer, Read Off a Graph}
Every equation above is really a question about where two graphs meet: the curve $y=b^{x}$ and a
horizontal line.
\begin{center}
\begin{minipage}[c]{0.42\linewidth}
\centering
\begin{tikzpicture}[scale=0.42]
  \lgrid{-3}{4}{-4}{9}
  \begin{scope}
    \clip (-3,-4) rectangle (4,9);
    \expgraph{1}{0.693147}{0}{0}{-3}{4}
  \end{scope}
  \solline{-3}{4}{4}
  \draw[redacc, dashed, very thick] (-3,-3)--(4,-3);
  \node[goldacc, font=\scriptsize, right] at (2.3,4.7) {$y=4$};
  \node[redacc, font=\scriptsize, right] at (2.3,-3.6) {$y=-3$};
  \fill[charcoal] (2,4) circle (3pt) node[above left, font=\scriptsize]{$(2,4)$};
\end{tikzpicture}
\end{minipage}
\begin{minipage}[c]{0.54\linewidth}
\small
\begin{itemize}[leftmargin=1.3em, itemsep=4pt]
  \item The gold line meets the curve at $x=$ \blank{1.2cm}, the solution of $2^{x}=4$ --- the same
        answer the algebra gave.
  \item This is also the point where $y=2^{x}-4$ crossed the $x$-axis in Lesson 6.2. Then you
        \emph{read} it off a picture; now you can \emph{solve} for it.
  \item The red line $y=-3$ never meets the curve, so $2^{x}=-3$ has \blank{2.6cm}.
  \item That is not an algebra failure --- it is the \textbf{range}. The range of $y=2^{x}$ is
        \blank{2.2cm}, and $-3$ is not in it.
\end{itemize}
\end{minipage}
\end{center}

\vspace{3pt}
So for $b^{x}=c$: if $c>0$ there is exactly \blank{1.6cm} solution; if $c\le0$ there are
\blank{1.6cm}. An equation of this shape never has two answers --- that $\pm$ belongs to the parabola
in the Hook, not here.

\vspace{5pt}
\begin{tcolorbox}[colback=sky, colframe=navy, boxrule=0.8pt, arc=1.5mm,
    left=2.5mm, right=2.5mm, top=1.5mm, bottom=1.5mm]
\textbf{A wall, on purpose.} Look at $2^{x}=5$. The line $y=5$ sits between $y=4$ and $y=8$, so it
\emph{does} cross the curve --- between $x=$ \blank{1.0cm} and $x=$ \blank{1.0cm}. The solution
exists. But $5$ is not a power of $2$, so Step 1 is impossible and today's method cannot reach it.

\vspace{2pt}
Keep these three. Every one of them has a solution; not one of them has a common base:
\[
  2^{x}=5 \qquad 132(0.85)^{t}=32 \;\;\text{(the coffee)} \qquad (0.9)^{t}=2.118 \;\;\text{(the soda)}
\]
That gap is what Lesson 6.4 walks straight up to, and what Unit 7 finally closes.
\end{tcolorbox}
\end{notesbox}

\begin{practicebox}
Solve each equation. Show the common-base step --- the answer alone is not the work.
\begin{enumerate}[leftmargin=1.7em, itemsep=9pt]
  \item $5^{x}=125$ \hfill $x=$ \blank{2.0cm}

  \item $\left(\frac{1}{2}\right)^{x}=16$ \hfill $x=$ \blank{2.0cm}

  \item $4^{\,x+1}=8$ \hfill $x=$ \blank{2.0cm}

  \item Explain why $3^{x}=0$ has no solution. Use the range and the asymptote, not algebra.
        \writelines{2}
\end{enumerate}
\end{practicebox}

\end{document}
